\documentclass[../dayan_kernel_manual.tex]{subfiles}

\begin{document}

\section{RPC 传输与处理}\label{sec:kernel-061}

RPC command 是一个 64-bit word：高 16 位为 opcode，低 48 位为参数。参数通常是
Device 地址；发送前必须确认高 16 位为 0。

\begin{center}
\code{command = (opcode << 48) | argument}
\end{center}

\subsection{mailbox 与 doorbell}

下列地址是相对 \code{CONFIG_CONTROL_BASE} 的 offset：

\begin{dayantable}{RPC 寄存器}{tab:kernel-065}{X[28,l] X[28,l] X[44,l]}
  寄存器 & offset & 用途 \\
  Host--Device mailbox & \code{0x1000 + hart*8} & Host 写入 command。 \\
  CLINT \code{msip} & \code{0x20000 + hart*4} & Host 写 1 提交；Device 写 0 确认。 \\
  Device--Host mailbox & \code{0x2020 + hart*8} & Device 写入 command。 \\
  \code{IRQ_STATUS} & \code{0x2000} & Device/Host 查询 outbound pending bit。 \\
  \code{IRQ_STATUS_SET} & \code{0x2008} & Device 置位本 hart outbound bit。 \\
  \code{IRQ_STATUS_CLR} & \code{0x200c} & 清除遗留 outbound 状态。 \\
  \code{IRQ_ENABLE_SET} & \code{0x2010} & 使能对应 Host IRQ。 \\
\end{dayantable}

两个方向使用相同的握手规则：发送方先写 mailbox，再置 doorbell；接收方处理完 command
后清除 doorbell。前一笔未确认时，发送方不能复用 mailbox。MMIO 两侧使用
\code{fence iorw,iorw} 保证 mailbox 与 doorbell 的顺序。

\subsection{Device 侧处理}

每个 hart 有 receive、send 和 profile 三个内部 proc，各使用独立的 4 KiB 栈；profile
proc 不参与普通 mailbox 收发，只在 trace FIFO 满时发起 \code{RPC_PROFILE_FLUSH}。

receive proc 由 \code{switch_context} 根据 \code{mip.MSIP} 直接选择，读取 Host command，调用
\code{rpc_receive(opcode,arg,response)}。handler 返回值的含义为：正数表示发送
response，0 表示无响应，负数表示请求已消费但处理失败。receive proc 同步发送 response，
忙等 Host 确认后清除 \code{msip}，整个事务期间不挂起；随后才切换到下一个上下文。

send proc 独占本 hart 的发送队列。队列为空时等待事件；取到 transaction 后，先等待
outbound bit 清零，写 mailbox 并置位，再忙等 Host 清零，发送期间不切出。最后执行 cleanup，标记事务
完成并唤醒调用者。root 自然退出使用 terminal transaction：调用者不再被唤醒，send
proc 在 Host 确认和 cleanup 之后直接回收已停止运行的 root，保证 transaction 所在栈
在握手完成前保持有效。

\code{rpc_call} 在调用 proc 的栈上保存 transaction，并给 packet、request、response
分别保留独占的 64-byte 对齐 slot。用户 Thread 等待 send proc 的完成事件，期间可以运行
其它 proc；系统 Thread 直接同步发送。内部调用者保证 opcode、结构体大小和参数有效，
\code{rpc_call} 不再重复检查这些约定。kill-all 摘除调用者的 event node，并取消仍在
send queue 的 transaction。发送期间不会切到 receive Thread，因此 kill 不会遇到
已发出但未完成握手的 transaction；已经返回硬件队列的调用者仍需消费 token 后才回收。

\subsection{Host 侧处理}

Service 启动独立后台 worker，调用 \code{npu_hal_irq_wait} 等待 outbound IRQ。HAL 在
超时时返回 \code{-ETIMEDOUT}；收到中断时返回 Cluster 编号和 IRQ status。具体后端可
等待事件，也可在 HAL 内部兼容轮询，Service 不区分两种实现。每个置位 Core 的处理顺序为：

\begin{enumerate}
  \item 读取 Device--Host mailbox；
  \item 若 opcode 和 argument 匹配该 hart 的待响应调用，记录 response；否则按
        第~\ref{sec:kernel-060}~节分派 Host 服务；
  \item 确认数据与状态已写回后，清除 IRQ bit；
  \item 唤醒等待 Host--Device response 的线程。
\end{enumerate}

Service 用每 Hart 的 launch agent 串行提交 CREATE，并由独立 RPC worker 接收回复
和退出通知。快速 kernel 的退出通知可能先于 launch agent 读完创建结果；两者通过
同一 mutex/condition 交接 \code{submit_done}，退出通知不会访问创建描述符。
Device 创建 proc 只使用静态 slot 和栈，不扩堆、不切出。

Host 发布算子模块的顺序是：先向每个已启动 hart 发送带模块范围的
\code{RPC_DCACHE_FLUSH}，全部完成后搬运新的 text/rodata，再把模块 entry 交给 launch
路径。每次 CREATE 在目标 Hart 执行 \code{fence.i} 后才允许 kernel 运行；模块复用时
不再失效 D-cache 或上传代码，也不需要单独的 I-cache RPC。

\subsection{cache ownership}

普通 Device RAM 不与 Host 硬件一致。RPC 用显式 clean/invalidate 交接数据：

\begin{dayantable}{RPC cache 顺序}{tab:kernel-066}{X[25,l] X[35,l] X[40,l]}
  对象 & 发布方 & 完成方 \\
  Device packet/request & 各占独立 64-B slot；Device clean 后置 outbound bit & Host 只读；处理完成后写 packet status 并清 bit。 \\
  Device response & 独占 64-B slot；Device 提交前 clean & Host 写入；Device 收到确认后 invalidate 并复制给调用者。 \\
  trap record & Device 复制到稳定槽并 clean & Host 读取并清 bit；确认前槽不可复用。 \\
  proc create 描述符 & Host 写入普通 Device RAM 后提交 & Device invalidate、复制启动数据、完整 flush 创建结果；Host 收到回复后读回并释放。 \\
  PS/kill 描述符 & Host 写入普通 Device RAM 后提交 & Device invalidate、填充结果并 clean；Host 收到 response 后读取。 \\
  D-cache flush 描述符 & Host 写入目标范围并提交 & Device invalidate 描述符和目标范围，回写 status 并 clean。 \\
  profile FIFO & Device 写 proc 初始状态与目标状态变化，满/STOP 时 clean & Host 读取后用 response 决定是否继续。 \\
\end{dayantable}

表中的 Host-owned 标量控制 descriptor 均从 64-byte 对齐地址开始，allocation 大小按
64-byte 向上取整，opcode 唯一确定结构体布局和长度。Device 在解析前 invalidate；
请求字段保持只读，只在响应前写结果并完整 flush。Host 收到回复并读回结果后释放。
提交前失败可以立即释放；一旦尝试发布命令却未收到匹配响应，该 allocation 不重新进入
分配池，只在 HAL 关闭时统一回收。

一致性窗口不需要上述普通 RAM 交接；Shared Named Resource 同步仍遵守
acquire/release 语义。

\subsection{并发与失败边界}

每个方向、每个 hart 只有一个硬件 mailbox，因此同一方向的事务由软件串行化；不同
hart 可以并行。Host 清除 doorbell 后，发送方可立即复用 mailbox 和相关缓冲，接收方
不得继续访问旧事务数据。

trap report 使用每 hart 单一稳定槽。前一报告未确认时再次异常、send proc 异常或
嵌套 trap 都会停机。Host 应把未知 opcode 和无效 packet 作为协议错误处理，并仍完成
必要的 doorbell 确认，避免通道永久阻塞。

\end{document}
